\section{Post-Iteration Processing}
\label{sec:misc-details}
%% ===================================================================

The per-pixel iteration kernels produce a buffer of raw escape-time
counts---integers that record how many iterations each pixel survived
before escaping.  These counts are not directly displayable: they must
be mapped to colors, and the resulting image must be antialiased to
suppress the severe spatial aliasing that arises at deep zooms (where
intricate fractal features routinely fall below pixel scale).  Without
careful post-processing, the visual output would be a noisy, jagged
approximation of the underlying mathematics.

This section covers the three main post-iteration stages:
antialiasing via GPU supersampling (\cref{subsec:antialiasing}), the
palette-based coloring system (\cref{subsec:coloring}), and
iteration-count statistics gathered via parallel GPU reduction.  The
subsequent section (\cref{sec:render-thread-pool}) describes the
threading architecture that orchestrates these stages alongside the UI.

% -----------------------------------------------------------------
\subsection{Antialiasing via supersampling}
\label{subsec:antialiasing}

\FractalShark{} supports GPU-side supersampled antialiasing (SSAA).
The iteration kernels render at a resolution that is a multiple of
the output size; a separate \emph{antialiasing kernel} then
down-samples the result to the final pixel grid.

The supersampling factor is a compile-time template parameter
\code{Antialiasing}, explicitly instantiated for $1\times$, $2\times$,
$3\times$, and $4\times$.  For a factor~$k$, each output pixel
corresponds to a $k\times k$ block in the high-resolution iteration
buffer \code{OutputIterMatrix}.  The kernel loops over the $k^2$
samples, maps each iteration count to a palette color
(\cref{subsec:coloring}), accumulates the RGB channels, and divides
by $k^2$ to produce an averaged \code{Color16} (16-bit per channel
RGBA) output.

\paragraph{Stream separation.}
Progressive frame extraction launches antialiasing and color readback
on \code{m\_DisplayStream}, separate from the iteration kernels'
\code{m\_ComputeStream}. Final frame extraction uses the compute stream.
GPU palette recoloring also uses the compute stream. GPU PNG saves run
coloring and antialiasing followed by encoding on that same stream
(\cref{subsec:png-output}).
The save caller waits for completion before releasing its snapshot.

\paragraph{Kernel template.}
\begin{lstlisting}
template<typename IterType,
         uint32_t Antialiasing,
         bool ScaledColor>
__global__ void antialiasing_kernel(
    const IterType *outputIterMatrix,
    uint32_t width, uint32_t height,
    AntialiasedColors outputColorMatrix,
    Palette pals,
    int colorWidth, int colorHeight,
    IterType numIterations,
    FractalShark::ColoringMode coloringMode);
\end{lstlisting}
The explicit coloring mode selects palette lookup or Basic grayscale
recoloring. The legacy \code{ScaledColor} template parameter is retained
but does not alter the current calculation. Samples whose original
iteration count reaches \code{numIterations} are left black.

\paragraph{Limitations.}
Brute-force SSAA supersamples every pixel uniformly, which is
wasteful: most of the image consists of smooth-gradient regions where
a single sample suffices, while visual aliasing concentrates along the
boundary of the Mandelbrot set.  An edge-aware or adaptive
supersampling strategy---sampling densely only near detected
edges---would achieve comparable image quality at a fraction of the
cost.  The current uniform approach was chosen for implementation
simplicity on the GPU (all threads perform identical work, avoiding
divergence), but replacing it with an adaptive scheme is a natural
future improvement.

% -----------------------------------------------------------------
\subsection{Coloring and palette system}
\label{subsec:coloring}

After the antialiasing kernel maps iteration counts to colors, the
mapping itself is controlled by the palette system implemented in
\code{FractalPalette}.

\paragraph{Palette types.}
Five palettes are available, enumerated by
\code{FractalPaletteType}:
\code{Basic} (placeholder),
\code{Default} (rainbow transitions R$\to$Y$\to$G$\to$C$\to$B$\to$M),
\code{Patriotic} (red/blue),
\code{Summer} (red/green/cyan/white), and
\code{Random} (procedurally generated).
Palette data is pre-generated on the CPU (using smooth color
transitions via \code{PalTransition}) and uploaded to the GPU as an
array of \code{Color16} values.

\paragraph{Bit-depth levels.}
Each palette is generated at six bit-depth levels, corresponding to
$2^5$, $2^6$, $2^8$, $2^{12}$, $2^{16}$, and $2^{20}$ colors
(32 to roughly one million).  The active depth is selected by
\code{m\_PaletteDepthIndex} (default: index~2, i.e.\ 256 colors)
and can be cycled at runtime.

\paragraph{Iteration-to-color mapping.}
On the GPU, the palette index for a given iteration count is computed
as
\[
  \mathrm{palIndex} \;=\;
  \bigl(\mathtt{numIters} \gg \mathtt{palette\_aux\_depth}\bigr)
    \bmod \mathtt{local\_palIters},
\]
Here, \code{numIters} denotes an escaped count after palette rotation and
saturation. The \code{palette\_aux\_depth} (range 0--16) controls the iteration
compression---higher values produce wider color bands---and
\code{local\_palIters} is the palette size at the active bit depth.
The modulo wrapping causes the palette to cycle when the iteration
count exceeds the palette length.

\paragraph{Palette rotation.}
CPU and GPU coloring add the palette rotation before the auxiliary-depth
shift and palette lookup, allowing a phase shift without recomputing the
iterations. Escape status is determined from the original count. Escaped
counts saturate at one less than the maximum permitted iteration count;
the GPU addition checks the available range before adding to avoid overflow.
Renderer initialization refreshes rotation, auxiliary depth, and the maximum
count even when the palette allocation is reused. GPU PNG saves supply the
settings captured by their save snapshot (\cref{subsec:png-output}).

\paragraph{Recoloring an existing frame.}
\code{EnqueuePaletteRecolor} applies a palette mutation and rebuilds colors
from the authoritative host iteration buffer without recalculating geometry.
For a GPU algorithm with CUDA enabled, the worker leases a renderer before
taking the calculation mutex, uploads logical rows with the padded host
stride, and calls the shared coloring and antialiasing kernel. A synchronous
readback updates the current host color buffer before the worker publishes
the frame through the existing presentation queue. Iteration counts remain
unchanged. Stream work is fenced on success, CUDA failure, and exception exits
before host storage or the renderer lease can be released. CPU algorithms and
instances that bypass the GPU retain the CPU recolorer.

Basic recoloring retains the CPU recolorer's grayscale formula, including
its iteration-dependent factor and 16-bit wrapping. The explicit mode leaves
normal GPU display and PNG palette lookup unchanged. Cancellation, superseding,
paced animation, and disabled-repaint placeholders retain the existing queue
behavior. Focused validation is described in \cref{subsec:png-validation}.

\paragraph{Alternative coloring approaches.}
The palette-indexed scheme above is simple and fast but has known
limitations: discrete iteration counts produce visible banding, and
the palette cycle length is fixed.  Several alternative techniques
exist in the fractal-rendering community:
\begin{itemize}
  \item \textbf{Smooth (continuous) iteration count.}  By incorporating
    the final orbit magnitude into a fractional escape count
    ($n_{\text{smooth}} = n - \log_2\!\log_2|z_n| + \text{const}$),
    banding is eliminated and color transitions become continuous.
  \item \textbf{Distance estimation.}  The analytic derivative
    $\partial z_n/\partial c$ (already computed for perturbation
    purposes) can yield an estimate of the distance from~$c$ to the
    Mandelbrot boundary.  Mapping distance to color emphasizes fine
    boundary structure and enables resolution-independent rendering of
    thin filaments.
  \item \textbf{Histogram equalization.}  Collecting a histogram of
    iteration counts across the image and mapping colors via the
    cumulative distribution produces uniform color utilization
    regardless of the iteration distribution---useful at deep zooms
    where most counts cluster in a narrow range.
  \item \textbf{Orbit-trap coloring.}  Measuring the closest approach
    of the orbit to a geometric shape (circle, cross, point) during
    iteration provides a distance metric unrelated to escape time,
    yielding distinctive visual effects.
\end{itemize}
\FractalShark{} currently implements only the palette-indexed method.
Integrating smooth iteration counts is feasible with minimal additional
complexity and would
pair naturally with the existing palette infrastructure; distance
estimation would leverage derivative data that the perturbation
pipeline already maintains.

% -----------------------------------------------------------------
\subsection{Iteration statistics via parallel reduction}
\label{subsec:reduction}

After the antialiasing kernel completes, a single reduction pass
extracts global statistics from the iteration buffer.  The
\code{max\_kernel} template reads every element of
\code{OutputIterMatrix} and computes three values simultaneously:
\begin{itemize}
  \item the minimum iteration count,
  \item the maximum iteration count,
  \item the sum of all iteration counts.
\end{itemize}
These are stored in a \code{ReductionResults} structure
(three \code{uint64\_t} fields: \code{Min}, \code{Max}, \code{Sum}).

The kernel uses shared-memory accumulators within each thread block
and atomic operations (\code{atomicMax}, \code{atomicMin},
\code{atomicAdd}) to merge block-level results into a single global
output.  It is launched with a $16\times16$ block grid and runs on
the same stream as the antialiasing kernel, so no additional
synchronization is required.  After the kernel completes, the host
copies back the \code{ReductionResults} via \code{cudaMemcpy} in
\code{ExtractItersAndColors}.

\paragraph{Current status.}
The reduction results are computed every frame but are not currently
used to drive any rendering decision (e.g.\ adaptive palette scaling
or histogram equalization).  They exist primarily for diagnostic
interest and as infrastructure that could support future coloring
enhancements.


% ===========================================================================
